\documentclass{article}
\usepackage[T1]{fontenc}
\usepackage[utf8]{inputenc}
\usepackage{ismir} % Remove the "submission" option for camera-ready version
\usepackage{amsmath,cite,url}
\usepackage{graphicx}
\usepackage{color}


% Title. Please use IEEE-compliant title case when specifying the title here,
% as it has implications for the copyright notice
% ------
\title{Transformer Based Symbolic Music Generation}

% Two addresses
% --------------
\twoauthors
   {Wongee Hong} {The Cooper Union \\ Albert Nerken School of Engineering}
   {Lennart Vetter} {The Cooper Union \\ Albert Nerken School of Engineering}

\def\authorname{W. Hong, and L. Vetter}

\sloppy % please retain sloppy command for improved formatting

\begin{document}
\maketitle

\begin{abstract}
This work addresses the limitations of recurrent architectures and small-scale datasets in 
symbolic music generation. Building upon previous attempts to steer music generation through Genetic Evolution
Algorithms (GEA) on LSTMs—which yielded suboptimal results in musical coherence—we propose a transition 
to a Transformer-based architecture. By aggregating three distinct datasets (VGMIDI, EMOPIA, and Pop1K7), 
we train a decoder-only Transformer on a large-scale corpus of approximately X million tokens. 
Our model employs a vocabulary of 1,350 unique tokens to represent melody and harmony. 
Preliminary results indicate that while the model learns basic structural dependencies, as 
evidenced by n-gram structural cosine similarity metrics, the transition from LSTM to Transformer 
significantly improves the stability of the generated musical sequences.
\end{abstract}

\section{Introduction}\label{sec:introduction}
Symbolic music generation has evolved from rule-based systems to sophisticated deep generative models capable 
of capturing complex polyphonic structures \cite{compoundword}. However, two primary challenges persist: 
the effective modeling of long-range dependencies in full-song compositions and the effective training 
on large scale datasets. 
Previous work \cite{vgmidi} attempted to use Long Short-Term Memory (LSTM) networks trained on the 
VGMIDI dataset, employing a Genetic Evolution Algorithm to steer output towards specific sentiments. 
However, these autoregressive models like LSTM and RNNs often suffered from "drift," where the lack of 
long-range attention led to a loss of structural integrity due to vanishing gradient signals 
during backpropagation in time for long sequences [CITATION NEEDED] thus yielding music pieces with incoherent 
strucutre.

In this project, we test the hypothesis that music generation is essentially a data-hungry task that 
requires high-capacity architectures capable of modeling global dependencies. We pivot from LSTMs to a 
Transformer architecture and expand our training data by a factor of 10, incorporating the EMOPIA and 
Pop1K7 datasets. We focus on the "quality-first" approach: establishing a baseline of high-quality music 
generation first and leaving the integration of sentiment steering up to further research.

\section{Related Work}

Symbolic Representations 
Standard MIDI-like representations treat every musical event (Note-On, Note-Off, Time-Shift) as a separate token. 
The REMI (REvamped MIDI) representation improved this by introducing metrical tokens like "Bar" and "Position" 
to provide a clear metric grid.

Transformer Models: The Music Transformer introduced relative attention to better capture periodicities. 
Subsequent models like MuseNet and Pop Music Transformer have applied these to diverse genres but often treat all 
token types equally in the embedding space. Recent surveys \cite{Ji} highlight the shift from 
RNNs to self-attention mechanisms, noting that Transformers better capture the periodic and 
hierarchical nature of music.

\section{Body}
To be done.

\subsection{Data}\label{subsec:data}
Datasets
The EMOPIA dataset \cite{emopia} provides a multi-modal (audio/MIDI) foundation for emotion-labeled pop piano. 
The VGMIDI dataset offers video game music \, while the Pop1K7 dataset [Hsiao et al., 2021] provides a large-scale repository of pop piano transcriptions from YouTube.
We use three different datasets in combination to train our model. We utilize the EMOPIA dataset which contains
1087 music clips from 387 sings [cite]. We further utilize the pop1k

\subsection{Proposed Method}\label{subsec:data_prep}
We adopt a vocabulary of $n=1,350$ tokens. The data preprocessing pipeline treats music modeling as a 
language modeling task. Musical events are encoded into features representing melody and harmony. 
Following the encoding scheme presented by \cite{emopia}, we use tokens representing pitch, 
duration, velocity or bar and position.

\subsection{Model Architecture}\label{subsec:model}
We utilize a decoder-only Transformer. The input to the model at time $t$ is a projected embedding of 
the music token. The architecture consists of multiple self-attention layers that 
compute:$$\text{Attention}(Q, K, V) = \text{softmax}\left(\frac{QK^T}{\sqrt{d_k}}\right)V$$

Where $d_k$ is the dimension of the key vectors. The model predicts the next music token word by 
employing multi-headed attention.

\section{Experimental Setup}\label{subsec:setup}

\subsection{Datasets}
We combine three primary datasets to overcome the "data-hungry" nature of Transformer models:

VGMIDI: 728 MIDI files (~202,030 tokens).

EMOPIA: 1,087 MIDI files (~480,335 tokens).

Pop1K7: 1,700 MIDI files (~1,577,835 tokens).

Total Corpus: 3,515 MIDI files, totaling approximately 2,260,200 tokens. We also experiment with transposition
and stretching of the dataset to further increase the training dataset size. The largest training dataset
used in the experiments totaled X M tokens.


\subsection{Training}\label{subsec:training}
The model was trained on the combined corpus using a cross-entropy loss function. We monitored 
training progress by inspecting generated sequences at various intervals to distinguish between 
chaotic "noise" and structured "music."

\section{Results}\label{sec:results}
To be done.

\subsection{Figures, Tables and Captions}
Include at least one table to display the results and \textbf{Graph of the Model Architecture}

Before training, the model produced random distributions of tokens 
(e.g., $t_120 v_56 d_whole_0 n_40 w_1599...$). After training, the output shifted 
toward valid musical syntax, correctly sequencing tempo ($t$), velocity ($v$), 
duration ($d$), and note ($n$) tokens (e.g., $t_120 v_56 d_eighth_0 v_60 d_quarter_0 n_79$).

To evaluate the structural consistency of the generated music, we measured the Structural Cosine Similarity across different n-gram levels:

\begin{table}
  \centering
  \begin{tabular}{|l|l|}
    \hline
    String value & Numeric value \\
    \hline
    Hello ISMIR  & \conferenceyear \\
    \hline
  \end{tabular}
  \caption{Table captions should be placed below the table.}
  \label{tab:example}
\end{table}

\section{Discussion and Error Analysis}
The transition from LSTM to Transformer addressed the "quality" bottleneck. However, the error analysis
reveals that the model still struggles with:

Metric Consistency: Occasionally skipping beats or losing the 4/4 time signature.
Harmonic Logic: While individual notes are correctly formatted, the global harmonic progression 
sometimes lacks a clear tonal center. The results confirm that the Transformer is less prone 
to the "unraveling" seen in LSTMs but requires even more data or more refined conditioning 
(like the emotion labels in EMOPIA) to reach professional-level composition.

\begin{thebibliography}{citations}
\bibitem{emopia} H.-T. Hung, J. Ching, S. Doh, N. Kim, J. Nam, and Y.-H. Yang, ``EMOPIA: A Multi-Modal Pop Piano Dataset For Emotion Recognition and Emotion-based Music Generation,'' \textit{arXiv preprint arXiv:2108.01374}, 2021.
\bibitem{vgmidi} L. N. Ferreira and J. Whitehead, ``Learning to generate music with sentiment,'' in \textit{Proc. 20th Int. Soc. Music Inf. Retrieval Conf. (ISMIR)}, 2019.
\bibitem{compoundword}W.-Y. Hsiao, J.-Y. Liu, Y.-C. Yeh, and Y.-H. Yang, ``Compound Word Transformer: Learning to Compose Full-Song Music over Dynamic Directed Hypergraphs,'' \textit{CoRR}, vol. abs/2101.02402, 2021. [Online]. Available: https://arxiv.org/abs/2101.02402
\bibitem{Ji2023} Ji, Shulei et al. “A Survey on Deep Learning for Symbolic Music Generation: Representations, Algorithms, Evaluations, and Challenges.” ACM Computing Surveys 56 (2023): 1 - 39.
\end{thebibliography}

\end{document}
